\documentclass[11pt]{article}
\usepackage{amsmath,amssymb,amsthm}
\newtheorem{proposition}{Proposition}
\begin{document}

\begin{proposition}[No general count-free closure]
Consider the finite-$M$ MSRC operator
\[
A_n=Q_M\otimes I_{d_n}+
\operatorname{blockdiag}(G_n(0)^T,\ldots,G_n(M)^T),
\]
where $G_n(k)$ contains background-dependent switching and coalescence rates.
For arbitrary $M$, switching rate, selection, branch durations, and initial
count distribution, the unconditional quartet probabilities do not admit a
general exact closure on a state space independent of $M$ based only on the
three topology probabilities.
\end{proposition}

\begin{proof}[Proof of the obstruction]
The topology value at an older boundary is a function $v(k,s)$ of the Moran
count $k$ and genealogy state $s$. In the structured generator, a pair in
background 1 coalesces at $1/\max(k,1)$ and a pair in background 0 at
$1/\max(M-k,1)$, while switching rates are proportional to $k/M$ and
$(M-k)/M$. Hence the generator maps a count-independent value function to a
function with distinct $k$-dependence whenever both coalescence and switching
are relevant.

Equivalently, the partition that identifies all count states is not strongly
lumpable: its rows have different aggregate rates to the topology-absorbing
states. Repeated generator application therefore produces additional
count-dependent functions, including inverse-count terms, so the three
topology masses are not a closed invariant coordinate system.

At a species split, daughter value functions are multiplied after being
conditioned on the same parental count. Thus replacing them by marginal means
loses their shared-count covariance. The explicit balanced example with
$M=10$ and initial distributions $\delta_5$ and
$(\delta_2+\delta_8)/2$ has equal mean count but different quartet vectors,
which rules out mean-frequency closure even for a neutral finite tree.
\end{proof}

The result does not exclude reductions at fixed frequency/count or under
special exchangeability symmetries. It also does not rule out matrix-free
application of the original count-indexed operator; that is an implementation
optimization, not a reduction in mathematical dimension.

\end{document}
